%%%% CAPÍTULO 2 - REVISÃO DA LITERATURA (OU REVISÃO BIBLIOGRÁFICA, ESTADO DA ARTE, ESTADO DO CONHECIMENTO)
%%
%% O autor deve registrar seu conhecimento sobre a literatura básica do assunto, discutindo e
% comentando a informação já publicada.
%% A revisão deve ser apresentada, preferencialmente, em ordem cronológica e por blocos de assunto,
% procurando mostrar a evolução do tema.
%% Título e rótulo de capítulo (rótulos não devem conter caracteres especiais, acentuados ou cedilha)
\chapter{Fundamentação Teórica}\label{cap:fundamentacaoTeorica}

Este capítulo apresenta os conceitos necessários para o entendimento da
proposta e dos experimentos. Primeiro, discute-se por que um sistema de
microsserviços é um sistema distribuído e que desafios isso traz. Em seguida,
distinguem-se as formas de comunicação síncrona e assíncrona, e apresenta-se o
\textit{Remote Procedure Call} (RPC), modelo sobre o qual o gRPC é construído.

\section{Microsserviços como sistemas distribuídos}

Em uma aplicação monolítica, os componentes compartilham o mesmo espaço de
endereçamento e se comunicam por chamadas de função diretas. Ao decompor a
aplicação em microsserviços, cada serviço passa a executar em seu próprio
processo, em geral em máquinas diferentes. O resultado é um sistema distribuído:
um conjunto de componentes em computadores interligados em rede que se
comunicam e coordenam suas ações apenas pela troca de mensagens
\cite{DistributedSystems:2013}.

Essa mudança tem duas consequências centrais para este trabalho. A primeira é a
latência: uma chamada que atravessa a rede é mais lenta que uma chamada local,
pois envolve transmissão, serialização dos dados e cópia entre \textit{buffers}
\cite{DistributedSystems:2013}. Assim, o custo de cada chamada deixa de ser
desprezível, e a forma como os serviços se comunicam passa a influenciar
diretamente o desempenho. A segunda são as falhas parciais: os componentes
podem falhar de forma independente enquanto os demais continuam em
funcionamento, e o cliente não consegue distinguir se a ausência de resposta
vem de uma falha no servidor ou na rede \cite{DistributedSystems:2013}. Por
isso, o sistema precisa ser projetado para detectar e tolerar falhas, e não
apenas evitá-las. A consistência dos dados é outro custo da distribuição, mas
está fora do escopo deste trabalho.

\section{Chamada de Procedimento Remoto (RPC)}

O RPC é um paradigma de \textit{middleware} que estende a chamada de
procedimento local para o ambiente distribuído. Ele permite que um cliente
invoque procedimentos de um servidor remoto com a mesma sintaxe de uma chamada
local, ocultando detalhes da distribuição, como a codificação e a decodificação
de parâmetros e resultados e a troca de mensagens \cite{DistributedSystems:2013}.

Essa transparência é viabilizada por um representante local do procedimento
remoto, chamado \textit{proxy} ou \textit{stub} no lado do cliente. A cada
chamada, ele empacota (\textit{marshalling}) a identificação da operação e seus
argumentos em uma mensagem de requisição, a envia ao servidor, bloqueia a
execução do cliente até a chegada da resposta e, por fim, desempacota
(\textit{unmarshalling}) os resultados ou exceções e os entrega ao código
chamador \cite{DistributedSystems:2013}. Esse trabalho representa um custo de
processamento a cada chamada, que se soma à latência da rede. Embora o RPC
oculte a distribuição do programador, ele não a elimina: a latência e as falhas
parciais discutidas anteriormente continuam presentes.

\section{gRPC}

% TODO: Protocol Buffers (contrato e serialização), HTTP/2 (multiplexação e
% streams), tipos de chamada (unária e streaming) e deadlines.
% Fonte sugerida: \cite{gRPC-Introduction:2026} e a documentação oficial.



